\documentclass[12pt,a4paper]{article}

% Packages
\usepackage[utf8]{inputenc}
\usepackage[T1]{fontenc}
\usepackage{geometry}
\usepackage{graphicx}
\usepackage{hyperref}
\usepackage{listings}
\usepackage{xcolor}
\usepackage{booktabs}
\usepackage{longtable}
\usepackage{enumitem}
\usepackage{float}
\usepackage{caption}
\usepackage{fancyhdr}
\usepackage{titlesec}
\usepackage{tocloft}
\usepackage{amsmath}

\geometry{margin=1in}

% Colors
\definecolor{codegreen}{rgb}{0,0.6,0}
\definecolor{codegray}{rgb}{0.5,0.5,0.5}
\definecolor{codepurple}{rgb}{0.58,0,0.82}
\definecolor{backcolour}{rgb}{0.95,0.95,0.95}
\definecolor{iitkblue}{rgb}{0.0,0.2,0.5}

% Code listing style
\lstdefinestyle{mystyle}{
    backgroundcolor=\color{backcolour},
    commentstyle=\color{codegreen},
    keywordstyle=\color{blue},
    numberstyle=\tiny\color{codegray},
    stringstyle=\color{codepurple},
    basicstyle=\ttfamily\small,
    breakatwhitespace=false,
    breaklines=true,
    captionpos=b,
    keepspaces=true,
    numbers=left,
    numbersep=5pt,
    showspaces=false,
    showstringspaces=false,
    showtabs=false,
    tabsize=2,
    frame=single,
    framerule=0.5pt,
    rulecolor=\color{codegray}
}
\lstset{style=mystyle}

% Header/Footer
\pagestyle{fancy}
\fancyhf{}
\fancyhead[L]{\small CSE 816: Software Production Engineering}
\fancyhead[R]{\small Final Project Report}
\fancyfoot[C]{\thepage}
\renewcommand{\headrulewidth}{0.4pt}

% Hyperref setup
\hypersetup{
    colorlinks=true,
    linkcolor=iitkblue,
    citecolor=iitkblue,
    urlcolor=iitkblue
}

% Title formatting
\titleformat{\section}{\Large\bfseries\color{iitkblue}}{\thesection}{1em}{}
\titleformat{\subsection}{\large\bfseries}{\thesubsection}{1em}{}
\titleformat{\subsubsection}{\normalsize\bfseries}{\thesubsubsection}{1em}{}

% ========================================================================
\begin{document}

% Title Page
\begin{titlepage}
    \centering
    \vspace*{1cm}

    {\Huge\bfseries\color{iitkblue} Fraud Detection DevOps Pipeline\par}
    \vspace{0.5cm}
    {\Large Real-Time Financial Fraud Detection with MLOps and DevSecOps\par}

    \vspace{2cm}

    {\large\textbf{CSE 816: Software Production Engineering}\par}
    {\large Final Project Report\par}

    \vspace{2cm}

    \begin{tabular}{ll}
        \toprule
        \textbf{Team Members} & \\
        \midrule
        Atul Soni & Roll No: 2025501 \\
        Raghav & \\
        \bottomrule
    \end{tabular}

    \vspace{3cm}

    {\large Department of Computer Science \& Engineering\par}

    \vspace{1cm}

    {\large\today\par}

    \vfill
\end{titlepage}

% Table of Contents
\tableofcontents
\newpage

% ========================================================================
\section{Introduction}

\subsection{Project Overview}

This project implements a \textbf{real-time financial fraud detection system} built with modern MLOps and DevSecOps practices. It combines an unsupervised machine learning model (Isolation Forest) served via a REST API with a complete CI/CD pipeline, container orchestration, centralized logging, and secrets management.

The project demonstrates a complete DevOps framework to automate the Software Development Life Cycle (SDLC) as required by the CSE 816 course, with a focus on the \textbf{MLOps and Finance} domain rather than a generic web application.

\subsection{Problem Statement}

Financial fraud causes billions of dollars in losses annually across the banking and e-commerce sectors. Manual fraud detection is slow, error-prone, and unable to scale with the volume of modern digital transactions. There is a clear need for an automated, scalable, real-time fraud scoring system backed by robust DevOps infrastructure for reliable production deployment.

\subsection{Objectives}

\begin{enumerate}
    \item Develop a machine learning-based fraud detection API capable of scoring financial transactions in real-time.
    \item Implement an end-to-end DevOps pipeline covering continuous integration, continuous deployment, monitoring, and security.
    \item Demonstrate industry-standard practices including containerization, orchestration, infrastructure-as-code, and centralized log management.
    \item Incorporate advanced features: HashiCorp Vault for secrets, Ansible Roles for modular configuration, and Kubernetes HPA for auto-scaling.
\end{enumerate}

% ========================================================================
\section{Tech Stack}

\begin{table}[H]
\centering
\caption{Technology Stack}
\begin{tabular}{@{}lll@{}}
\toprule
\textbf{Layer} & \textbf{Technology} & \textbf{Purpose} \\
\midrule
Application & Python 3.11, FastAPI, Pydantic & REST API with request validation \\
Machine Learning & scikit-learn (Isolation Forest) & Anomaly detection model \\
Serialization & joblib, NumPy, Pandas & Model persistence and data processing \\
Version Control & Git, GitHub & Source code management \\
CI/CD & Jenkins, GitHub Webhooks & Automated pipeline \\
Containerization & Docker, Docker Compose & Application packaging \\
Config Management & Ansible (with Roles) & Infrastructure-as-code \\
Orchestration & Kubernetes, HPA & Auto-scaling and orchestration \\
Secrets Management & HashiCorp Vault & Secure credential storage \\
Monitoring \& Logging & ELK Stack & Centralized log management \\
Testing & pytest, FastAPI TestClient & Automated unit tests \\
\bottomrule
\end{tabular}
\end{table}

% ========================================================================
\section{Project Structure}

\begin{lstlisting}[language={},caption={Repository Structure}]
fraud-detection-devops/
|-- app/                        # Application layer
|   |-- main.py                 # FastAPI app (/predict, /health)
|   |-- train.py                # Isolation Forest model training
|   |-- simulate.py             # Live transaction simulator
|   |-- test_main.py            # Unit tests (pytest)
|   |-- Dockerfile              # Container image definition
|   +-- requirements.txt        # Python dependencies
|-- jenkins/                    # CI/CD pipeline
|   +-- Jenkinsfile             # 7-stage declarative pipeline
|-- docker/                     # Docker Compose stack
|   +-- docker-compose.yml      # 5 services: API, ELK, Vault
|-- ansible/                    # Configuration management
|   |-- site.yml                # Main playbook
|   |-- inventory.ini           # Host inventory
|   +-- roles/
|       |-- app/main.yml        # Deploy fraud-api container
|       |-- monitoring/main.yml # Deploy ELK stack
|       +-- k8s/main.yml        # Apply K8s manifests
|-- k8s/                        # Kubernetes manifests
|   |-- deployment.yaml         # RollingUpdate + health probes
|   |-- service.yaml            # LoadBalancer service
|   +-- hpa.yaml                # Horizontal Pod Autoscaler
|-- elk/                        # Log pipeline config
|   +-- logstash.conf           # JSON log parsing
+-- vault/                      # Secrets
    +-- vault-policy.hcl        # Vault access control policy
\end{lstlisting}

% ========================================================================
\section{Application Design}

\subsection{Machine Learning Model}

\textbf{Algorithm:} Isolation Forest (unsupervised anomaly detection from scikit-learn).

The Isolation Forest algorithm isolates anomalies by randomly selecting a feature and then randomly selecting a split value between the maximum and minimum values of the selected feature. Anomalous observations require fewer random splits to be isolated, resulting in shorter path lengths in the tree structure.

The model is trained on the \textbf{Kaggle Sparkov Fraud Detection Dataset} (\textasciitilde1.3 million transactions) using the following engineered features:

\begin{table}[H]
\centering
\caption{Feature Engineering}
\begin{tabular}{@{}lll@{}}
\toprule
\textbf{Feature} & \textbf{Type} & \textbf{Derivation} \\
\midrule
\texttt{amt} & float & Transaction amount (raw) \\
\texttt{category\_enc} & int & Merchant category (Label Encoded) \\
\texttt{hour} & int & Hour extracted from transaction timestamp \\
\texttt{distance} & float & Euclidean distance (cardholder $\leftrightarrow$ merchant) in km \\
\texttt{is\_fraud} & int & Ground truth fraud label \\
\bottomrule
\end{tabular}
\end{table}

\textbf{Model Hyperparameters:}
\begin{itemize}
    \item \texttt{n\_estimators = 100} (number of isolation trees)
    \item \texttt{contamination = 0.01} (expected fraction of anomalies)
    \item \texttt{random\_state = 42} (reproducibility)
\end{itemize}

The trained model is serialized using \texttt{joblib} as \texttt{fraud\_model.pkl}.

\begin{lstlisting}[language=Python,caption={Model Training (train.py)}]
from sklearn.ensemble import IsolationForest
from sklearn.preprocessing import LabelEncoder
import joblib, pandas as pd, numpy as np

df = pd.read_csv("fraudTrain.csv")
df["hour"] = pd.to_datetime(df["trans_date_trans_time"]).dt.hour
df["distance"] = np.sqrt(
    (df["lat"] - df["merch_lat"])**2 +
    (df["long"] - df["merch_long"])**2
) * 111  # approx km

le = LabelEncoder()
df["category_enc"] = le.fit_transform(df["category"])

FEATURE_COLS = ["amt", "category_enc", "hour", "distance", "is_fraud"]
X = df[FEATURE_COLS].dropna()

model = IsolationForest(n_estimators=100, contamination=0.01,
                        random_state=42)
model.fit(X[FEATURE_COLS])
joblib.dump(model, "fraud_model.pkl")
\end{lstlisting}

\subsection{API Design}

The application is built with \textbf{FastAPI}, a modern Python web framework known for high performance and automatic OpenAPI documentation.

\begin{table}[H]
\centering
\caption{API Endpoints}
\begin{tabular}{@{}llll@{}}
\toprule
\textbf{Endpoint} & \textbf{Method} & \textbf{Description} & \textbf{Response} \\
\midrule
\texttt{/} & GET & Root status & \texttt{\{"message": "..."\}} \\
\texttt{/predict} & POST & Fraud prediction & \texttt{\{"is\_fraud": bool, "fraud\_score": float\}} \\
\texttt{/health} & GET & Health check & \texttt{\{"status": "ok"\}} \\
\bottomrule
\end{tabular}
\end{table}

\textbf{Request Schema (\texttt{/predict}):}

\begin{lstlisting}[language={}]
{
    "amount": 4500.0,
    "merchant_category": "online_retail",
    "hour_of_day": 2,
    "distance_from_home_km": 1200.0,
    "is_foreign_transaction": 1
}
\end{lstlisting}

\textbf{Response Schema:}

\begin{lstlisting}[language={}]
{
    "is_fraud": true,
    "fraud_score": -0.5,
    "merchant_category": "online_retail",
    "amount": 4500.0
}
\end{lstlisting}

The \texttt{merchant\_category} field maps to encoded values: \texttt{grocery}~(0), \texttt{online\_retail}~(1), \texttt{gas}~(2), \texttt{entertainment}~(3), \texttt{travel}~(4), \texttt{other}~(5).

\subsection{Transaction Simulator}

\texttt{simulate.py} generates continuous random transactions for testing and log generation:

\begin{itemize}
    \item \textbf{Normal transactions (90\%):} Low amounts (\$5--\$300), daytime hours, short distances, domestic.
    \item \textbf{Suspicious transactions (10\%):} High amounts (\$800--\$5000), late-night hours, long distances, foreign.
\end{itemize}

% ========================================================================
\section{Version Control (Git \& GitHub)}

The project uses \textbf{Git} for version control and \textbf{GitHub} as the remote repository. Key practices implemented:

\begin{itemize}
    \item All source code, configuration files, and infrastructure definitions are version-controlled.
    \item \texttt{.gitignore} excludes sensitive and generated files: \texttt{\_\_pycache\_\_}, \texttt{*.pkl}, \texttt{*.csv}, \texttt{.env}, \texttt{app.log}.
    \item GitHub Webhooks trigger the Jenkins CI/CD pipeline automatically on every push to the \texttt{main} branch.
    \item Incremental updates to the repository trigger the full automated pipeline as required.
\end{itemize}

% ========================================================================
\section{CI/CD Pipeline (Jenkins)}

\subsection{Pipeline Overview}

The Jenkins pipeline is defined as a \textbf{Declarative Jenkinsfile} and triggers automatically on every push to the \texttt{main} branch via \texttt{githubPush()} webhook trigger.

\subsection{Pipeline Stages}

The pipeline consists of \textbf{7 automated stages}:

\begin{figure}[H]
\centering
\begin{verbatim}
  +----------+    +---------------+    +---------------------+
  | Checkout |-->>| Fetch Secrets |-->>| Install Dependencies|
  +----------+    | (from Vault)  |    +---------------------+
                  +---------------+              |
                                                 v
  +---------------------+    +------------------+    +-----------+
  | Deploy with Ansible |<<--| Push to DockerHub|<<--| Run Tests |
  +---------------------+    +------------------+    +-----------+
                                                          |
                                                          v
                                              +---------------------+
                                              | Build Docker Image  |
                                              +---------------------+
\end{verbatim}
\caption{Jenkins Pipeline Flow}
\end{figure}

\begin{table}[H]
\centering
\caption{Jenkins Pipeline Stages}
\begin{tabular}{@{}clp{8cm}@{}}
\toprule
\textbf{Stage} & \textbf{Name} & \textbf{Description} \\
\midrule
1 & Checkout & Clones the \texttt{main} branch from GitHub \\
2 & Fetch Secrets & Retrieves credentials from HashiCorp Vault \\
3 & Install Dependencies & Runs \texttt{pip install -r requirements.txt} \\
4 & Run Tests & Executes \texttt{pytest} with verbose output \\
5 & Build Docker Image & Builds container image tagged with build number \\
6 & Push to Docker Hub & Authenticates and pushes image to Docker Hub \\
7 & Deploy with Ansible & Runs Ansible playbook to deploy the new image \\
\bottomrule
\end{tabular}
\end{table}

\begin{lstlisting}[language={},caption={Jenkinsfile (Abbreviated)}]
pipeline {
    agent any
    triggers { githubPush() }

    environment {
        DOCKERHUB_CREDENTIALS = credentials('dockerhub-creds')
        IMAGE_NAME = "fraud-detection-api"
        IMAGE_TAG = "${BUILD_NUMBER}"
        VAULT_ADDR = "http://vault:8200"
    }

    stages {
        stage('Checkout') {
            steps {
                git branch: 'main', url: "https://github.com/..."
            }
        }
        stage('Fetch Secrets from Vault') {
            steps {
                sh '''
                  export VAULT_TOKEN=$(cat /run/secrets/vault_token)
                  export DB_PASSWORD=$(vault kv get -field=password
                    secret/fraud-app/db)
                '''
            }
        }
        stage('Run Tests') {
            steps {
                sh 'cd app && python -m pytest test_main.py -v'
            }
        }
        stage('Build Docker Image') {
            steps {
                sh "docker build -t ${DOCKERHUB_CREDENTIALS_USR}/
                    ${IMAGE_NAME}:${IMAGE_TAG} ./app"
            }
        }
        stage('Push to Docker Hub') {
            steps {
                sh '''
                  echo $DOCKERHUB_CREDENTIALS_PSW | docker login
                    -u $DOCKERHUB_CREDENTIALS_USR --password-stdin
                  docker push ${DOCKERHUB_CREDENTIALS_USR}/
                    ${IMAGE_NAME}:${IMAGE_TAG}
                '''
            }
        }
        stage('Deploy with Ansible') {
            steps {
                sh "ansible-playbook -i ansible/inventory.ini
                    ansible/site.yml
                    -e image_tag=${IMAGE_TAG}
                    -e dockerhub_user=${DOCKERHUB_CREDENTIALS_USR}"
            }
        }
    }
    post {
        success { echo "Pipeline succeeded." }
        failure { echo "Pipeline failed." }
    }
}
\end{lstlisting}

\subsection{Automated Trigger Mechanism}

The pipeline is triggered via \textbf{GitHub Webhook} using the \texttt{githubPush()} trigger. When a developer pushes code to the \texttt{main} branch:

\begin{enumerate}
    \item GitHub sends a POST webhook to the Jenkins server.
    \item Jenkins receives the webhook and triggers the pipeline.
    \item The full 7-stage pipeline executes automatically.
    \item On success, the new version is deployed; on failure, the team is notified.
\end{enumerate}

This fulfills the requirement that \textit{``incremental updates to the Git repository should trigger automated processes''}.

% ========================================================================
\section{Containerization (Docker)}

\subsection{Dockerfile}

The application is containerized using a lightweight \texttt{python:3.11-slim} base image for minimal footprint:

\begin{lstlisting}[language={},caption={Dockerfile}]
FROM python:3.11-slim
WORKDIR /app
COPY requirements.txt .
RUN pip install --no-cache-dir -r requirements.txt
COPY . .
EXPOSE 8000
CMD ["uvicorn", "main:app", "--host", "0.0.0.0", "--port", "8000"]
\end{lstlisting}

\subsection{Docker Compose}

The \texttt{docker-compose.yml} orchestrates \textbf{5 services} in a single stack:

\begin{table}[H]
\centering
\caption{Docker Compose Services}
\begin{tabular}{@{}lllp{5cm}@{}}
\toprule
\textbf{Service} & \textbf{Image} & \textbf{Port} & \textbf{Purpose} \\
\midrule
\texttt{fraud-api} & Built from Dockerfile & 8000 & Fraud detection API \\
\texttt{elasticsearch} & elasticsearch:8.12.0 & 9200 & Log storage engine \\
\texttt{logstash} & logstash:8.12.0 & 5044 & Log ingestion \& parsing \\
\texttt{kibana} & kibana:8.12.0 & 5601 & Log visualization \\
\texttt{vault} & vault:1.16 & 8200 & Secrets management \\
\bottomrule
\end{tabular}
\end{table}

\textbf{Key Design Decisions:}
\begin{itemize}
    \item \textbf{Shared volumes:} \texttt{app\_logs} volume is shared between \texttt{fraud-api} (write) and \texttt{logstash} (read-only) for log forwarding.
    \item \textbf{Persistent storage:} \texttt{es\_data} volume ensures Elasticsearch data survives container restarts.
    \item \textbf{Service dependencies:} \texttt{depends\_on} ensures correct startup order.
    \item \textbf{Environment variables:} \texttt{LOG\_FILE} and \texttt{ENV} are injected via environment configuration.
\end{itemize}

% ========================================================================
\section{Configuration Management (Ansible)}

\subsection{Playbook Architecture}

Ansible is used for \textbf{automated deployment} and \textbf{infrastructure provisioning}. The playbook follows a \textbf{modular role-based architecture} as encouraged by the evaluation criteria.

\begin{lstlisting}[language={},caption={site.yml --- Main Playbook}]
---
- name: Deploy Fraud Detection Pipeline
  hosts: all
  become: yes
  vars:
    image_tag: "latest"
  roles:
    - app
    - monitoring
    - k8s
\end{lstlisting}

\subsection{Role: \texttt{app} --- Application Deployment}

Handles deploying the fraud detection API container:

\begin{enumerate}
    \item \textbf{Pull} the latest Docker image from Docker Hub (tagged with Jenkins build number).
    \item \textbf{Stop} any existing \texttt{fraud-api} container.
    \item \textbf{Start} a new container with \texttt{restart\_policy: always} on port 8000.
\end{enumerate}

\begin{lstlisting}[language={},caption={Role: app/main.yml}]
---
- name: Pull latest fraud-api Docker image
  docker_image:
    name: "{{ dockerhub_user }}/fraud-detection-api:{{ image_tag }}"
    source: pull

- name: Stop existing container
  docker_container:
    name: fraud-api
    state: absent

- name: Start fraud-api container
  docker_container:
    name: fraud-api
    image: "{{ dockerhub_user }}/fraud-detection-api:{{ image_tag }}"
    state: started
    restart_policy: always
    ports:
      - "8000:8000"
\end{lstlisting}

\subsection{Role: \texttt{monitoring} --- ELK Stack Deployment}

Deploys the complete ELK monitoring stack:

\begin{enumerate}
    \item \textbf{Elasticsearch} --- single-node mode on port 9200.
    \item \textbf{Logstash} --- with pipeline config for JSON log parsing on port 5044.
    \item \textbf{Kibana} --- connected to Elasticsearch on port 5601.
\end{enumerate}

\subsection{Role: \texttt{k8s} --- Kubernetes Deployment}

Applies all Kubernetes manifests:

\begin{enumerate}
    \item \texttt{deployment.yaml} --- Application deployment with RollingUpdate.
    \item \texttt{service.yaml} --- LoadBalancer service.
    \item \texttt{hpa.yaml} --- Horizontal Pod Autoscaler.
\end{enumerate}

% ========================================================================
\section{Container Orchestration (Kubernetes)}

\subsection{Deployment Configuration}

\begin{table}[H]
\centering
\caption{Kubernetes Deployment Specification}
\begin{tabular}{@{}ll@{}}
\toprule
\textbf{Parameter} & \textbf{Value} \\
\midrule
Replicas & 2 (minimum) \\
Update Strategy & RollingUpdate \\
Max Surge & 1 \\
Max Unavailable & 0 (zero-downtime) \\
CPU Request & 100m \\
CPU Limit & 500m \\
Memory Request & 128Mi \\
Memory Limit & 512Mi \\
\bottomrule
\end{tabular}
\end{table}

\subsection{Health Probes}

\begin{itemize}
    \item \textbf{Readiness Probe:} \texttt{GET /health} --- initial delay 10s, period 5s. Ensures traffic is only sent to ready pods.
    \item \textbf{Liveness Probe:} \texttt{GET /health} --- initial delay 15s, period 10s. Automatically restarts unhealthy pods.
\end{itemize}

\subsection{Service}

The Kubernetes Service exposes the API externally:

\begin{itemize}
    \item \textbf{Type:} LoadBalancer
    \item \textbf{Port mapping:} External port 80 $\rightarrow$ Container port 8000
    \item Provides a stable external IP/DNS for accessing the API.
\end{itemize}

\subsection{Horizontal Pod Autoscaler (HPA)}

\begin{table}[H]
\centering
\caption{HPA Configuration}
\begin{tabular}{@{}ll@{}}
\toprule
\textbf{Parameter} & \textbf{Value} \\
\midrule
Min Replicas & 2 \\
Max Replicas & 10 \\
Scaling Metric & CPU Utilization \\
Target Utilization & 60\% \\
\bottomrule
\end{tabular}
\end{table}

The HPA automatically scales the number of pods between 2 and 10 based on CPU utilization, ensuring the system handles traffic spikes without manual intervention.

\subsection{Zero-Downtime Deployment (Live Patching)}

The combination of \texttt{RollingUpdate} strategy with \texttt{maxUnavailable: 0} ensures:

\begin{enumerate}
    \item New pods are created with the updated image.
    \item Readiness probes confirm the new pods are healthy.
    \item Traffic is gradually shifted to new pods.
    \item Old pods are terminated only after new ones are ready.
\end{enumerate}

This implements \textbf{live patching} --- the application is updated without any downtime. Upon refreshing the application after deployment, the new changes are visible seamlessly.

% ========================================================================
\section{Monitoring \& Logging (ELK Stack)}

\subsection{Logging Pipeline Architecture}

\begin{figure}[H]
\centering
\begin{verbatim}
  +------------+    +-----------+    +---------------+    +---------+
  | FastAPI    |-->>| JSON Log  |-->>|   Logstash    |-->>| Elastic |
  | Application|    | File      |    | (Parse+Index) |    | Search  |
  +------------+    +-----------+    +---------------+    +----+----+
                                                              |
                                                              v
                                                        +---------+
                                                        | Kibana  |
                                                        | Dashboard|
                                                        +---------+
\end{verbatim}
\caption{ELK Logging Pipeline}
\end{figure}

\subsection{How It Works}

\begin{enumerate}
    \item \textbf{Application Logging:} Every \texttt{/predict} request generates a structured JSON log entry containing: timestamp, amount, merchant category, hour, distance, foreign transaction flag, fraud score, and fraud status.

    \item \textbf{Log Collection:} Logstash reads the JSON log file from a shared Docker volume (\texttt{app\_logs}).

    \item \textbf{Log Parsing:} Logstash parses the JSON, extracts the \texttt{@timestamp} field (ISO8601), and adds a \texttt{service: fraud-api} tag.

    \item \textbf{Log Storage:} Parsed logs are indexed in Elasticsearch with daily indices: \texttt{fraud-logs-YYYY.MM.dd}.

    \item \textbf{Visualization:} Kibana provides search, filtering, and dashboards for real-time fraud monitoring.
\end{enumerate}

\subsection{Logstash Configuration}

\begin{lstlisting}[language={},caption={logstash.conf}]
input {
  file {
    path => "/app/logs/app.log"
    start_position => "beginning"
    sincedb_path => "/dev/null"
    codec => "json"
  }
}

filter {
  date {
    match => ["timestamp", "ISO8601"]
    target => "@timestamp"
  }
  mutate {
    add_field => { "service" => "fraud-api" }
  }
}

output {
  elasticsearch {
    hosts => ["http://elasticsearch:9200"]
    index => "fraud-logs-%{+YYYY.MM.dd}"
  }
  stdout { codec => rubydebug }
}
\end{lstlisting}

\subsection{Sample Log Entry}

\begin{lstlisting}[language={}]
{
    "timestamp": "2026-05-12T18:05:00.123456",
    "amount": 4500.0,
    "merchant_category": "online_retail",
    "hour_of_day": 2,
    "distance_from_home_km": 1200.0,
    "is_foreign_transaction": 1,
    "fraud_score": -0.5,
    "is_fraud": true
}
\end{lstlisting}

This fulfills the requirement that \textit{``application logs must feed into the ELK Stack, and the Kibana dashboard should visualize these logs''}.

% ========================================================================
\section{Secrets Management (HashiCorp Vault)}

\subsection{Configuration}

HashiCorp Vault provides \textbf{secure, centralized secrets management}, ensuring sensitive credentials are never stored in code or configuration files.

\begin{itemize}
    \item Vault runs in dev mode via Docker Compose on port 8200.
    \item Root token: \texttt{root} (dev mode only; production would use proper authentication).
    \item IPC\_LOCK capability is enabled for memory protection.
\end{itemize}

\subsection{Access Policy}

The \texttt{vault-policy.hcl} defines a \textbf{least-privilege access policy}:

\begin{lstlisting}[language={},caption={vault-policy.hcl}]
path "secret/fraud-app/*" {
  capabilities = ["read", "list"]
}

path "secret/fraud-app/db" {
  capabilities = ["read"]
}

path "secret/fraud-app/dockerhub" {
  capabilities = ["read"]
}
\end{lstlisting}

\subsection{Integration Points}

\begin{itemize}
    \item \textbf{Jenkins Pipeline:} Stage 2 (``Fetch Secrets from Vault'') retrieves database and Docker Hub credentials at runtime.
    \item \textbf{Secret Paths:}
    \begin{itemize}
        \item \texttt{secret/fraud-app/db} --- Database credentials.
        \item \texttt{secret/fraud-app/dockerhub} --- Docker Hub credentials.
    \end{itemize}
    \item Credentials are fetched dynamically, never hardcoded in the codebase.
\end{itemize}

% ========================================================================
\section{Testing}

\subsection{Test Framework}

\begin{itemize}
    \item \textbf{Framework:} pytest with FastAPI \texttt{TestClient}.
    \item \textbf{Mocking:} \texttt{unittest.mock} patches \texttt{joblib.load} to inject a mock model, so tests run without the actual \texttt{.pkl} file.
    \item \textbf{HTTP Client:} httpx-based TestClient for synchronous API testing.
\end{itemize}

\subsection{Test Coverage}

\begin{table}[H]
\centering
\caption{Test Cases}
\begin{tabular}{@{}llp{6.5cm}@{}}
\toprule
\textbf{Test Class} & \textbf{Test Case} & \textbf{Validates} \\
\midrule
\texttt{TestHealthEndpoint} & \texttt{test\_health\_returns\_ok} & \texttt{/health} returns 200 with \texttt{\{"status":"ok"\}} \\
\midrule
\texttt{TestRootEndpoint} & \texttt{test\_root\_returns\_message} & \texttt{/} returns API running message \\
\midrule
\texttt{TestPredictEndpoint} & \texttt{test\_predict\_returns\_200} & Valid input returns 200 \\
 & \texttt{test\_predict\_response\_fields} & Response has all required fields \\
 & \texttt{test\_predict\_fraud\_detection} & Prediction $-1$ maps to fraud \\
 & \texttt{test\_predict\_normal\_transaction} & Prediction $1$ maps to normal \\
 & \texttt{test\_predict\_unknown\_category} & Unknown categories default to ``other'' \\
 & \texttt{test\_predict\_correct\_features} & Feature array shape is $(1, 5)$ \\
 & \texttt{test\_predict\_invalid\_payload} & Invalid types return 422 \\
 & \texttt{test\_predict\_empty\_payload} & Empty body returns 422 \\
\bottomrule
\end{tabular}
\end{table}

Tests are executed automatically in the Jenkins pipeline (Stage 4) on every code push.

% ========================================================================
\section{Architecture Diagram}

\begin{figure}[H]
\centering
\begin{verbatim}
                  +------------------------------------------+
                  |              DEVELOPER                    |
                  |          git push to GitHub               |
                  +-------------------+----------------------+
                                      |
                                      v
                  +------------------------------------------+
                  |         JENKINS CI/CD PIPELINE            |
                  |                                          |
                  |  Checkout >> Vault >> Install >> Test >>  |
                  |  Docker Build >> Push >> Ansible Deploy   |
                  +--------+-----------------------+---------+
                           |                       |
              +------------v-----------+   +-------v--------+
              |      DOCKER HUB       |   | HASHICORP VAULT |
              |   (Image Registry)    |   | (Secrets Store) |
              +------------+-----------+   +----------------+
                           |
              +------------v-------------------------------+
              |          ANSIBLE DEPLOYMENT                 |
              |   Role: app | Role: monitoring | Role: k8s |
              +-----+---------------+---------------+------+
                    |               |               |
                    v               v               v
            +-------+------+ +-----+-------+ +-----+--------+
            |  FRAUD API   | |  ELK STACK  | |  KUBERNETES  |
            |  (Docker)    | |             | |              |
            |  Port: 8000  | | ES: 9200   | | 2-10 Pods    |
            |              | | Kibana:5601 | | HPA (60%CPU) |
            |  /predict    | | Logstash    | | RollingUpdate|
            |  /health     | |             | | LoadBalancer |
            +---------+----+ +------+------+ +--------------+
                      |             ^
                      | JSON Logs   |
                      +-------------+
\end{verbatim}
\caption{Complete System Architecture}
\end{figure}

% ========================================================================
\section{Evaluation Criteria Mapping}

This section maps the project implementation to the evaluation criteria specified in the project guidelines.

\subsection{Working Project (20 Marks)}

\begin{table}[H]
\centering
\caption{Mandatory Requirements Mapping}
\begin{tabular}{@{}p{7cm}lp{4cm}@{}}
\toprule
\textbf{Requirement} & \textbf{Status} & \textbf{Implementation} \\
\midrule
Git push triggers automated processes & \checkmark & GitHub Webhook + \texttt{githubPush()} \\
Jenkins fetches and builds updated code & \checkmark & Stages 1, 3, 5 \\
Running automated tests & \checkmark & Stage 4 (pytest) \\
Pushing Docker images to Docker Hub & \checkmark & Stage 6 \\
Deploying images to target system & \checkmark & Stage 7 (Ansible) \\
New changes visible on refresh & \checkmark & RollingUpdate + zero downtime \\
Logs feed into ELK Stack & \checkmark & JSON $\rightarrow$ Logstash $\rightarrow$ ES \\
Kibana dashboard visualizes logs & \checkmark & \texttt{fraud-logs-*} index \\
\bottomrule
\end{tabular}
\end{table}

\subsection{Advanced Features (3 Marks)}

\begin{table}[H]
\centering
\caption{Advanced Features}
\begin{tabular}{@{}lp{9cm}@{}}
\toprule
\textbf{Feature} & \textbf{Implementation} \\
\midrule
\textbf{Vault} & HashiCorp Vault with least-privilege policy; Jenkins fetches secrets at runtime \\
\textbf{Ansible Roles} & 3 modular roles: \texttt{app}, \texttt{monitoring}, \texttt{k8s} \\
\textbf{Kubernetes HPA} & Auto-scales 2--10 pods at 60\% CPU utilization \\
\textbf{Live Patching} & RollingUpdate with \texttt{maxUnavailable: 0} for zero-downtime deploys \\
\bottomrule
\end{tabular}
\end{table}

\subsection{Innovation (2 Marks)}

\begin{table}[H]
\centering
\caption{Innovative Features}
\begin{tabular}{@{}lp{9cm}@{}}
\toprule
\textbf{Innovation} & \textbf{Description} \\
\midrule
\textbf{Domain: MLOps + Finance} & Real-world fraud detection instead of generic web app \\
\textbf{Isolation Forest} & Unsupervised anomaly detection for fraud scoring \\
\textbf{Real-Time Scoring API} & Sub-second fraud predictions via FastAPI \\
\textbf{Transaction Simulator} & Automated load generator with realistic fraud patterns \\
\textbf{Structured JSON Logging} & Every prediction logged with full context for analytics \\
\bottomrule
\end{tabular}
\end{table}

% ========================================================================
\section{How to Run}

\subsection{Prerequisites}

\begin{itemize}
    \item Python 3.11+
    \item Docker Desktop
    \item Kaggle dataset: \texttt{fraudTrain.csv} placed in \texttt{app/}
\end{itemize}

\subsection{Local Development}

\begin{lstlisting}[language=bash,caption={Local Setup Commands}]
# Train the model
cd app && python train.py

# Run unit tests
python -m pytest test_main.py -v

# Start the API server
uvicorn main:app --host 0.0.0.0 --port 8000

# Run the transaction simulator
python simulate.py
\end{lstlisting}

\subsection{Docker Compose (Full Stack)}

\begin{lstlisting}[language=bash,caption={Docker Compose Setup}]
cd docker
docker-compose up -d
\end{lstlisting}

Access points after startup:
\begin{itemize}
    \item API: \url{http://localhost:8000}
    \item Kibana: \url{http://localhost:5601}
    \item Elasticsearch: \url{http://localhost:9200}
    \item Vault: \url{http://localhost:8200}
\end{itemize}

\subsection{Kubernetes Deployment}

\begin{lstlisting}[language=bash,caption={Kubernetes Deployment}]
kubectl apply -f k8s/deployment.yaml
kubectl apply -f k8s/service.yaml
kubectl apply -f k8s/hpa.yaml
\end{lstlisting}

% ========================================================================
\section{Conclusion}

This project demonstrates a complete end-to-end MLOps pipeline for real-time financial fraud detection, covering the full lifecycle from model training to production deployment with monitoring. The system integrates:

\begin{itemize}
    \item \textbf{Machine Learning} --- Isolation Forest for unsupervised anomaly detection.
    \item \textbf{CI/CD Automation} --- Jenkins with 7 automated stages triggered by GitHub webhooks.
    \item \textbf{Containerization} --- Docker with a 5-service Compose stack.
    \item \textbf{Configuration Management} --- Ansible with 3 modular roles.
    \item \textbf{Orchestration} --- Kubernetes with HPA, health probes, and zero-downtime deployments.
    \item \textbf{Monitoring} --- ELK Stack with structured log pipeline and Kibana visualization.
    \item \textbf{Security} --- HashiCorp Vault for secrets management with least-privilege policies.
\end{itemize}

The project aligns with real-world DevOps practices in the \textbf{Finance and MLOps} domain, going beyond generic web applications to demonstrate how DevOps methodologies can be applied to production machine learning systems.

% ========================================================================

\section*{References}

\begin{enumerate}
    \item FastAPI Documentation: \url{https://fastapi.tiangolo.com/}
    \item scikit-learn Isolation Forest: \url{https://scikit-learn.org/stable/modules/generated/sklearn.ensemble.IsolationForest.html}
    \item Jenkins Pipeline Syntax: \url{https://www.jenkins.io/doc/book/pipeline/syntax/}
    \item Docker Compose Documentation: \url{https://docs.docker.com/compose/}
    \item Ansible Roles: \url{https://docs.ansible.com/ansible/latest/playbook_guide/playbooks_reuse_roles.html}
    \item Kubernetes HPA: \url{https://kubernetes.io/docs/tasks/run-application/horizontal-pod-autoscale/}
    \item ELK Stack: \url{https://www.elastic.co/elastic-stack}
    \item HashiCorp Vault: \url{https://www.vaultproject.io/docs}
    \item Kaggle Fraud Detection Dataset: \url{https://www.kaggle.com/datasets/kartik2112/fraud-detection}
\end{enumerate}

\end{document}
